% ============================================================
%   Chapter 3  --  Project Budget
% ============================================================

\chapter{Project Budget}\label{ch:budget}

This chapter estimates the cost of CALA-SafeRL \emph{as if it had been carried out as a
professional engineering contract}, rather than as an unpaid academic exercise. The estimate
builds directly on the planning of Chapter~\ref{ch:planning}: the labour cost follows the
role-and-effort distribution of the human-resources plan (Table~\ref{tab:hr-roles}), and the
contingency reserve follows the risk register (Table~\ref{tab:risks}).
Section~\ref{sec:budget-assumptions} fixes the cost model and its assumptions, the subsequent
sections cost labour, equipment, software, energy and travel, and
Section~\ref{sec:budget-total} aggregates them into the total budget.

\section{Cost model and assumptions}\label{sec:budget-assumptions}

The estimate rests on the following assumptions, chosen to be defensible for a project executed
in Spain during the 2025\,--\,2026 academic year:

\begin{itemize}
  \item \textbf{Labour} is charged per role at the rates of Table~\ref{tab:hr-roles}
        (\(40\)\,€/h Project Manager, \(35\)\,€/h RL/ML Engineer, \(30\)\,€/h Software Engineer,
        \(25\)\,€/h QA/Test Engineer, \(22\)\,€/h Technical Writer), reflecting Spanish market
        rates for a junior-to-mid professional in each speciality.
  \item \textbf{Equipment} is \emph{pre-owned} and charged by linear amortisation,
        \(\text{cost}\times(\text{months used}/\text{useful life})\times\text{dedication}\),
        with a \(60\)-month useful life for the workstation, an \(11\)-month project span and an
        \(80\%\) dedication factor.
  \item \textbf{Electricity} is priced at \(0.18\)\,€/kWh, a representative Spanish domestic
        tariff for the period.
  \item \textbf{Indirect costs} (workspace, internet, administration) are taken as \(15\%\) of
        direct costs. A \(5\%\) \textbf{contingency} reserve covers the risks of
        Table~\ref{tab:risks}.
  \item \textbf{Value-added tax} (IVA) is applied at the Spanish general rate of \(21\%\) on the
        base budget.
\end{itemize}

\section{Labour costs}\label{sec:budget-labor}

Labour is by far the dominant cost, as is typical for a research-and-development project whose
main input is engineering time. Table~\ref{tab:labor-internal} costs the author's own effort by
role, and Table~\ref{tab:labor-support} costs the support team (Director, Coordinator and tribunal),
whose time is a real cost of the project even though it is not borne by the author.

\begin{table}[H]
  \centering
  \caption{Internal labour cost (project team). Hours taken from the human-resources plan
           (Table~\ref{tab:hr-roles}).}
  \label{tab:labor-internal}
  \small
  \begin{tabular}{@{}lccr@{}}
    \toprule
    \textbf{Role} & \textbf{Hours} & \textbf{Rate (€/h)} & \textbf{Cost (€)} \\
    \midrule
    Project Manager   & 45  & 40 & 1\,800.00 \\
    RL/ML Engineer    & 256 & 35 & 8\,960.00 \\
    Software Engineer & 134 & 30 & 4\,020.00 \\
    QA/Test Engineer  & 45  & 25 & 1\,125.00 \\
    Technical Writer  & 70  & 22 & 1\,540.00 \\
    \midrule
    \textbf{Subtotal} & \textbf{550} & & \textbf{17\,445.00} \\
    \bottomrule
  \end{tabular}
\end{table}

\begin{table}[H]
  \centering
  \caption{Support-team labour cost.}
  \label{tab:labor-support}
  \small
  \begin{tabular}{@{}lccr@{}}
    \toprule
    \textbf{Member} & \textbf{Hours} & \textbf{Rate (€/h)} & \textbf{Cost (€)} \\
    \midrule
    Project Director (supervisor) & 30 & 60 & 1\,800.00 \\
    TFG Coordinator               & 5  & 50 & 250.00 \\
    Defense tribunal (3 members)  & 12 & 50 & 600.00 \\
    \midrule
    \textbf{Subtotal} & \textbf{47} & & \textbf{2\,650.00} \\
    \bottomrule
  \end{tabular}
\end{table}

The combined labour cost is therefore \(17\,445.00 + 2\,650.00 = 20\,095.00\)~€.

\section{Equipment and materials}\label{sec:budget-equipment}

The project ran on a single pre-owned desktop workstation (the machine described in
Section~\ref{sec:req-constraints} and used for the experiments of
Chapter~\ref{ch:experimentation}). Table~\ref{tab:equipment} charges only the fraction of its
useful life consumed by the project. No consumables of note were required, and, as discussed
in Section~\ref{sec:budget-software}, the entire scientific software stack is open-source, so
its licensing cost is zero.

\begin{table}[H]
  \centering
  \caption{Amortised equipment cost.
           Amortisation \(=\text{cost}\times(\text{months used}/\text{useful life})\times\text{dedication}\).}
  \label{tab:equipment}
  \small
  \begin{tabular}{@{}lccccr@{}}
    \toprule
    \textbf{Item} & \textbf{Cost (€)} & \textbf{Life (mo)} & \textbf{Used (mo)} & \textbf{Ded.} & \textbf{Amort. (€)} \\
    \midrule
    Workstation (i9-9900K, RTX 2080 Ti, 32\,GB) & 2\,000.00 & 60 & 11 & 80\% & 293.33 \\
    Monitor and peripherals                     & 350.00    & 72 & 11 & 80\% & 42.78 \\
    \midrule
    \textbf{Subtotal} & & & & & \textbf{336.11} \\
    \bottomrule
  \end{tabular}
\end{table}

\section{Software and AI services}\label{sec:budget-software}

A deliberate design decision was to build the framework entirely on open-source software, which
keeps the licensing cost at zero and the work fully reproducible (Table~\ref{tab:software}). The
only paid software is the AI coding-and-writing assistant whose use is disclosed in
Appendix~\ref{app:ai-usage}. The local Gemma model (served through Ollama) and the free tier of
Gemini incur no per-call cost, with their compute accounted for under energy
(Section~\ref{sec:budget-energy}).

\begin{table}[H]
  \centering
  \caption{Software and AI-service costs. The scientific stack is fully open-source.}
  \label{tab:software}
  \small
  \begin{tabular}{@{}p{5.8cm}p{5.0cm}r@{}}
    \toprule
    \textbf{Item} & \textbf{Licence / basis} & \textbf{Cost (€)} \\
    \midrule
    CARLA, Gymnasium, PyTorch, NumPy, Matplotlib, Streamlit, Plotly, SQLite, Ruff, pytest
      & MIT / BSD-3 / Apache-2.0 / PSF / public domain & 0.00 \\
    AI coding/writing assistant (\(\approx 20\)\,€/mo \(\times\) 8 mo)
      & Commercial subscription (Appendix~\ref{app:ai-usage}) & 160.00 \\
    Local LLM (Ollama + Gemma), Gemini free tier
      & Local / free tier & 0.00 \\
    \midrule
    \textbf{Subtotal} & & \textbf{160.00} \\
    \bottomrule
  \end{tabular}
\end{table}

\section{Energy consumption}\label{sec:budget-energy}

Reinforcement-learning training is energy-intensive: each configuration requires hundreds of
episodes of CARLA simulation with the policy on the GPU. Table~\ref{tab:energy} estimates the
electricity consumed, separating heavy GPU training from lighter development and analysis.

\begin{table}[H]
  \centering
  \caption{Energy consumption. Electricity price: \(0.18\)\,€/kWh.}
  \label{tab:energy}
  \small
  \begin{tabular}{@{}lccr@{}}
    \toprule
    \textbf{Use} & \textbf{Hours} & \textbf{Power (kW)} & \textbf{Energy (kWh)} \\
    \midrule
    GPU training (CARLA + policy) & 350 & 0.45 & 157.5 \\
    Development and analysis      & 250 & 0.20 & 50.0 \\
    \midrule
    \textbf{Total} & & & \textbf{207.5} \\
    \bottomrule
  \end{tabular}
\end{table}

At \(0.18\)\,€/kWh, the \(207.5\,\mathrm{kWh}\) consumed amount to \(37.35\)~€. While small in
absolute terms, the figure is reported for completeness and because it scales linearly with the
number of training runs, which is the project's main computational cost driver.

\section{Travel and accommodation}\label{sec:budget-travel}

The project was developed locally and supervised through a combination of on-campus and online
meetings, so no travel or accommodation was required. This line is therefore \(0.00\)~€.

\section{Total budget}\label{sec:budget-total}

Table~\ref{tab:budget-summary} aggregates the previous sections. Indirect costs are applied at
\(15\%\) of direct costs, the contingency reserve at \(5\%\) of the direct-plus-indirect subtotal,
and IVA at \(21\%\) of the base budget. The resulting total cost of the project is
\(\mathbf{30\,139.73}\)~€.

\begin{table}[H]
  \centering
  \caption{Total project budget.}
  \label{tab:budget-summary}
  \small
  \begin{tabular}{@{}lr@{}}
    \toprule
    \textbf{Concept} & \textbf{Cost (€)} \\
    \midrule
    Labour, project team (Table~\ref{tab:labor-internal})  & 17\,445.00 \\
    Labour, support team (Table~\ref{tab:labor-support})   & 2\,650.00 \\
    Equipment, amortised (Table~\ref{tab:equipment})          & 336.11 \\
    Software and AI services (Table~\ref{tab:software})       & 160.00 \\
    Energy (Table~\ref{tab:energy})                           & 37.35 \\
    Travel and accommodation                                  & 0.00 \\
    \midrule
    \textbf{Direct costs}                                     & \textbf{20\,628.46} \\
    Indirect costs / overhead (15\%)                          & 3\,094.27 \\
    Contingency reserve (5\%)                                 & 1\,186.14 \\
    \midrule
    \textbf{Base budget (before tax)}                         & \textbf{24\,908.87} \\
    VAT (IVA, 21\%)                                            & 5\,230.86 \\
    \midrule
    \textbf{Total project budget}                             & \textbf{30\,139.73} \\
    \bottomrule
  \end{tabular}
\end{table}

Two observations close the chapter. First, labour accounts for \(20\,095\)~€ of the
\(20\,628\)~€ of direct cost (roughly \(97\%\)), which confirms that, like most
research-and-development efforts, the project's value is overwhelmingly in engineering time rather
than in hardware or licences. Second, the open-source software stack contributes a licensing cost
of exactly zero. Had the framework been built on commercial simulation or RL tooling, this line
alone could have rivalled the hardware cost. The low marginal cost of an additional experiment
(essentially only electricity, \(\approx 0.18\)\,€/kWh) is precisely what makes a reproducible,
open Safe-RL framework such as CALA-SafeRL economically attractive as a research platform.

\clearpage
